%% ====================================================================
%% Pipelines
%% ====================================================================

\section[Pipelines]{Running pipelines with Snakemake}
\begin{frame}{What is a pipeline?}
A pipeline is a series of analysis steps, run one after the other
\begin{itemize}
    \item Each step's output is the next step's input
    \item The same steps run for every sample
\end{itemize}

\bigskip
\diagram{pipeline-steps}

\bigskip
A loop in a job script is already a small pipeline.
A \emph{workflow manager} does the bookkeeping for you:
which step comes next, what is done, what failed
\end{frame}

\begin{frame}{One job script or a Snakemake pipeline?}
\diagram{jobscript-vs-snakemake}
\medskip

\footnotesize
\begin{cmdtablex}{lLL}
                  & \textbf{One job script}                & \textbf{A Snakemake pipeline}             \\
    \midrule
    A step fails  & Start everything again                 & Fix it: only failed and missing steps run \\
    Resources     & Booked for the biggest step, all along & Each step books what it needs             \\
\end{cmdtablex}
\end{frame}

\begin{frame}{We use Snakemake}
\begin{itemize}
    \item All the lab's pipelines are Snakemake pipelines
    \item It uses what you learned today:
          Conda for the tools, Slurm for the jobs
    \item A step is a \emph{rule}: input files in, output files out
    \item Snakemake checks which files exist, and runs only what is missing
\end{itemize}

\bigskip
You run the lab's pipelines: you do not have to write them
\end{frame}

\begin{frame}{An example: \texttt{readqc}, the pipeline you run today}
\diagram{readqc-dag}

\bigskip
\begin{itemize}
    \item \cmd{fastp} trims the reads, \cmd{seqkit_stats} counts them,
          \cmd{merge_stats} puts all samples in one table
    \item Snakemake works out the order from the files
    \item Every step runs as its own Slurm job, many at the same time
    \item In \cmd{HPC_intro/readqc}, laid out like the lab's real pipelines
\end{itemize}
\end{frame}

\begin{frame}[fragile]{A step is a \emph{rule}: from input files to output files}
\begin{lstlisting}[style=scriptsize]
rule seqkit_stats:
    input:
        "results/{sample}/01_trim/{sample}_trimmed_R1.fastq.gz",
        "results/{sample}/01_trim/{sample}_trimmed_R2.fastq.gz",
    output:
        "results/{sample}/02_stats/{sample}_stats.tsv",
    conda:
        "../envs/qc.yaml"
    log:
        "logs/seqkit_stats/{sample}.log",
    shell:
        "seqkit stats -a -T {input} > {output} 2> {log}"
\end{lstlisting}
\hint{Simplified from \cmd{readqc/workflow/rules/qc.smk}}

\medskip
\footnotesize
\begin{itemize}
    \item \cmd{input} and \cmd{output}: the files the step needs and the files it makes
    \item \cmd{{sample}} is filled in for every sample in your sample sheet
    \item \cmd{shell}: the command; \cmd{{input}}, \cmd{{output}} and \cmd{{log}}
          are filled in from above
\end{itemize}
\end{frame}

\begin{frame}{Rules are linked by their files}
\diagram{rule-files}

\bigskip
\begin{itemize}
    \item Snakemake looks at which rule makes each input file, and works out the order
    \item You never tell it the order yourself
    \item A file that is already there is not made again,
          unless its input or its settings changed
\end{itemize}
\end{frame}

% Tree comments: a short gap instead of dotted leaders
\renewcommand*\DTcomment[1]{\hspace{2em}{\DTstylecomment #1}}
\begin{frame}{Inside a pipeline directory}
\DTsetlength{0.2em}{1em}{0.2em}{0.4pt}{1.6pt}
\renewcommand*\DTstyle{\ttfamily}
\renewcommand*\DTstylecomment{\sffamily\color{darkgray}}
\dirtree{%
.1 readqc/.
.2 README.md\DTcomment{how to install and run it}.
.2 \alert{config/}\DTcomment{\alert{yours}}.
.3 \alert{config.yaml}\DTcomment{the settings}.
.3 \alert{sample\_table.csv}\DTcomment{your samples}.
.2 workflow/\DTcomment{the pipeline: leave it as it is}.
.2 results/, logs/\DTcomment{appear when you run it}.
}
\end{frame}

\begin{frame}{The sample sheet: which reads belong to which sample}
\cmd{config/sample_table.csv}

\medskip
\footnotesize
\begin{cmdtable}{lll}
    \cmd{sample_name} & \cmd{fastq_r1}                 & \cmd{fastq_r2}                 \\
    \midrule
    \cmd{sample_A}    & \cmd{../sample_A_R1.fastq.gz}  & \cmd{../sample_A_R2.fastq.gz}  \\
    \cmd{sample_B}    & \cmd{../sample_B_R1.fastq.gz}  & \cmd{../sample_B_R2.fastq.gz}  \\
    \cmd{sample_C}    & \cmd{../sample_C_R1.fastq.gz}  & \cmd{../sample_C_R2.fastq.gz}  \\
\end{cmdtable}
\end{frame}

\begin{frame}[fragile]{Running a pipeline on the HPC}
In \cmd{readqc}, on the login node (fine: Snakemake only submits jobs)
\begin{lstlisting}[style=scriptsize]
conda activate snakemake       # create it once: exercise 5.1
# 1. sample sheet: which reads belong to which sample
nano config/sample_table.csv

# 2. dry run: what would run? Nothing runs yet
snakemake -j unlimited -n \
    --workflow-profile workflow/profiles/hpc_slurm/
#  -n                  dry run: only show the steps
#  -j unlimited        as many Slurm jobs at once as needed
#  --workflow-profile  the directory with the Slurm settings
#  \ at the end        the command goes on on the next line

# 3. run it: the same command without -n
#    first run: a few minutes while Snakemake installs the tools
\end{lstlisting}
\hint{Run it in \cmd{screen}: it keeps going until the last job is done}
\smallskip
\hint{On your laptop, use \cmd{--workflow-profile workflow/profiles/local/}}
\end{frame}

\begin{frame}{Exercises: a Snakemake pipeline}
\begin{exlist}{5_pipelines.md}
    \ex{5.1}{Set up the pipeline}{1}
    \ex{5.2}{Dry run}{1}
    \ex{5.3}{Run it}{1}
    \ex{5.4}{Run it again}{2}
    \ex{5.5}{Change a setting}{3}
\end{exlist}
\end{frame}

\begin{frame}{What you learned today}
\begin{itemize}
    \item Move around, work with files, \cmd{find}, \cmd{grep}, loops and scripts
    \item Save your work with Git and GitHub
    \item Install tools in Conda environments
    \item Log in to the HPC, run jobs with \cmd{srun} and \cmd{sbatch}, and run a pipeline
\end{itemize}

\bigskip
Those 200 FASTQ files? Check them with \cmd{md5sum}, keep your script in Git,
install the assembler with Conda, and run it with \cmd{sbatch}
\end{frame}
